\section{Desarrollo del taller}

\subsection{Diagnóstico inicial del caso}

\scenario{1}{Análisis dirigido de MediSalud · fuente literal}

La revisión del caso identifica tres procesos de máxima criticidad. Primero, la atención médica y el registro HCE, porque cualquier demora o pérdida de información afecta la continuidad clínica. Segundo, la prescripción y dispensación, donde un error puede causar daño directo. Tercero, la admisión y facturación, que constituyen la entrada del paciente y sostienen la operación financiera.

Los perfiles más expuestos son el médico, que necesita información completa durante una consulta; el paciente, que agenda, consulta resultados y usa telemedicina; y el personal de admisión, que resuelve disponibilidad, cobertura y cobro. Enfermería y farmacia también requieren sincronización oportuna con la HCE. Gerencia y calidad necesitan una vista agregada, pero no sustituyen la experiencia de los perfiles operativos.

\begin{table}[H]
\centering
\caption{Matriz de análisis inicial}
\begin{tabularx}{\textwidth}{p{3.4cm}Y}
\toprule
Pregunta guía & Respuesta consolidada \\
\midrule
Procesos críticos & Atención y HCE; prescripción y dispensación; agenda, admisión y facturación. \\
Usuarios afectados & Médicos, pacientes y admisión como grupos principales; enfermería y farmacia por dependencia clínica. \\
Evidencia disponible & Incidentes con fecha, módulo, descripción, rol y sede; contexto tecnológico y requisitos no funcionales. \\
Evidencia faltante & Resultados de tarea, duraciones, encuestas, escenarios de red/dispositivo y criterios uniformes de decisión. \\
Decisión metodológica & Preservar incidentes y complementar la observación con datos simulados identificados. \\
\bottomrule
\end{tabularx}
\end{table}

\subsection{Clasificación según ISO/IEC 25022}

\scenario{2}{Clasificación de 3.000 incidentes · fuente literal}

La clasificación automatizada asignó 1.853 casos a efectividad, 347 a eficiencia, 102 a satisfacción, 359 a libertad de riesgo y 339 a cobertura del contexto. Efectividad concentra el 61,77\% de los registros; sin embargo, la frecuencia no equivale a prioridad. Los 359 casos de libertad de riesgo requieren evaluación inmediata por la gravedad de sus consecuencias potenciales.

\begin{figure}[H]
\centering
\begin{tikzpicture}[x=0.0038cm,y=0.72cm]
  \foreach \label/\value/\y/\barcolor in {Efectividad/1853/5/mediteal,Eficiencia/347/4/mediblue,Satisfacción/102/3/medigold,Libertad de riesgo/359/2/medicoral,Cobertura del contexto/339/1/medigrey}{
    \fill[\barcolor!18] (0,\y-.22) rectangle (2000,\y+.22);
    \fill[\barcolor] (0,\y-.22) rectangle (\value,\y+.22);
    \node[anchor=east,font=\small] at (-25,\y) {\label};
    \node[anchor=west,font=\small\bfseries,text=\barcolor] at (\value+25,\y) {\value};
  }
  \draw[medigrey!50] (0,.45) -- (0,5.55);
\end{tikzpicture}
\caption{Distribución de incidentes por característica de calidad en uso.}
\label{fig:classification}
\end{figure}

\begin{table}[H]
\centering
\caption{Criterio y ejemplo de clasificación}
\begin{tabularx}{\textwidth}{p{3.2cm}YY}
\toprule
Característica & Pregunta de decisión & Ejemplo de efecto \\
\midrule
Efectividad & ¿La tarea no se completa o su resultado es incorrecto? & Una cita no queda confirmada pese a finalizar el formulario. \\
Eficiencia & ¿El resultado se alcanza con tiempo o pasos excesivos? & El guardado de la nota supera ocho segundos. \\
Satisfacción & ¿La experiencia reduce confianza, agrado o disposición de uso? & El paciente abandona una encuesta o reporta confusión. \\
Libertad de riesgo & ¿Existe daño clínico, financiero o de privacidad? & Una dosis incorrecta, factura duplicada o dato de otro paciente. \\
Cobertura del contexto & ¿El fallo depende de sede, dispositivo o conectividad? & Una teleconsulta cae bajo una red doméstica limitada. \\
\bottomrule
\end{tabularx}
\end{table}

La justificación se almacena junto a cada fila procesada. Esta práctica es necesaria porque una descripción puede admitir más de una lectura. Una receta que no se guarda afecta efectividad; una receta guardada con dosis errónea se clasifica principalmente como libertad de riesgo. La característica principal expresa el efecto dominante, aunque el análisis pueda reconocer impactos secundarios.

\evidence{La salida completa se encuentra en \texttt{data/processed/}\allowbreak\texttt{incidentes\_clasificados.csv}; el resumen de conteos está en \texttt{evidencias/resultados/}\allowbreak\texttt{clasificacion\_resumen.json}.}

\subsection{Atributos, tareas y prioridades}

\scenario{4--5}{Atributos UTC y mapeo tarea--característica · reconstrucción declarada}

Se seleccionaron tareas observables que cubren los perfiles del caso y pueden producir evidencia. La prioridad combina impacto clínico, financiero, operativo y reputacional con la urgencia sugerida por los incidentes.

\begin{longtable}{p{2.0cm}p{3.2cm}p{3.7cm}p{2.7cm}p{1.6cm}}
\caption{Catálogo de tareas, contextos y atributos}\label{tab:utc}\\
\toprule
Usuario & Tarea & Contexto & Atributo observable & Prioridad \\
\midrule
\endfirsthead
\toprule
Usuario & Tarea & Contexto & Atributo observable & Prioridad \\
\midrule
\endhead
Médico & Guardar nota HCE & Consulta, hora pico & P90 de duración & Crítica \\
Médico & Revisar alergias & Antes de prescribir & Disponibilidad correcta & Crítica \\
Paciente & Agendar cita & Móvil y 4G & Tasa de éxito & Alta \\
Paciente & Realizar teleconsulta & WiFi doméstico & Tasa de caídas & Alta \\
Admisión & Emitir factura & Seguro y cierre mensual & Tasa de error & Crítica \\
Farmacia & Validar receta & Turno nocturno & Recetas correctas & Crítica \\
Enfermería & Registrar signos & Tableta y red hospitalaria & Sincronización & Alta \\
Médico & Cargar imagen & Dispositivo clínico & P90 de carga & Alta \\
Paciente & Responder encuesta & Aplicación móvil & Abandono y CSAT & Media \\
Gerencia & Revisar KPI & Consolidado multisede & Cobertura de datos & Alta \\
\bottomrule
\end{longtable}

Las tareas críticas no son necesariamente las más frecuentes. La prescripción se prioriza por severidad; la HCE por su relación con la decisión clínica; y la facturación por el riesgo de cobros incorrectos. El portal de citas y la telemedicina tienen prioridad alta porque son puntos de acceso directo del paciente y dependen de contextos heterogéneos.

\subsection{Diseño de métricas}

\scenario{6}{Fichas de medida y criterios de decisión · reconstrucción declarada}

Cada ficha incluye código, nombre, característica, objetivo, fórmula, unidad, fuente, periodicidad, meta y semáforo. El valor no se interpreta sin su meta: 96,52\% puede parecer alto para una función general, pero es rojo cuando el resultado restante representa recetas incorrectas.

Las fórmulas principales son:

\begin{align}
\text{Éxito}(T) &= \frac{\text{tareas exitosas}}{\text{tareas iniciadas}}\times 100, \\
\text{Error}(T) &= \frac{\text{tareas no exitosas}}{\text{tareas iniciadas}}\times 100, \\
\text{Abandono} &= \frac{\text{encuestas no completadas}}{\text{encuestas iniciadas}}\times 100, \\
\text{CSAT} &= \frac{\sum_{i=1}^{n} \text{puntuación}_{i}}{n}, \\
\text{Variabilidad} &= \sigma\left(\overline{t}_{\text{sede }1},\ldots,\overline{t}_{\text{sede }5}\right).
\end{align}

El percentil 90 se define como el menor valor que contiene al 90\% de las observaciones ordenadas. Se eligió para HCE e imagenología porque hace visible la experiencia lenta del extremo superior sin depender de un único máximo atípico.

\begin{longtable}{p{1.6cm}p{3.6cm}p{2.6cm}p{2.0cm}p{2.1cm}}
\caption{Catálogo resumido de métricas}\label{tab:catalog}\\
\toprule
Código & Métrica & Característica & Unidad & Meta \\
\midrule
\endfirsthead
\toprule
Código & Métrica & Característica & Unidad & Meta \\
\midrule
\endhead
M-EFI-01 & P90 registro HCE & Eficiencia & segundos & $\leq 8$ \\
M-EFE-01 & Éxito de agendamiento & Efectividad & porcentaje & $\geq 95$ \\
M-RIE-01 & Errores de facturación & Libertad de riesgo & porcentaje & $<1$ \\
M-SAT-01 & Abandono de encuesta móvil & Satisfacción & porcentaje & $<10$ \\
M-CC-01 & Variabilidad HCE entre sedes & Cobertura del contexto & segundos & $\leq 2$ \\
M-RIE-02 & Incidentes de privacidad & Libertad de riesgo & incidentes & $=0$ \\
M-SAT-02 & CSAT promedio & Satisfacción & escala 1--5 & $\geq 4$ \\
M-EFE-02 & Recetas correctas & Efectividad & porcentaje & $=100$ \\
M-EFI-02 & P90 carga de imagen & Eficiencia & segundos & $\leq 10$ \\
M-CC-02 & Caídas de teleconsulta & Cobertura del contexto & porcentaje & $<5$ \\
\bottomrule
\end{longtable}

\subsection{Datos, automatización y KPI}

\scenario{7--9}{Obtención, procesamiento y visualización · reconstrucción declarada}

La simulación instrumenta cada flujo con un identificador de auditoría y un escenario. En operación normal se registra la finalización esperada; en modo defectuoso se activa una regla controlada como demora, pérdida de guardado, pasos adicionales, doble facturación, dosis incompatible o desconexión. La probabilidad y la semilla permiten repetir la distribución sin convertir el fallo en una vulnerabilidad real.

El pipeline analítico ejecuta cuatro responsabilidades separadas:

\begin{enumerate}
  \item \textbf{Clasificación:} carga los incidentes, normaliza campos y asigna característica y justificación.
  \item \textbf{Simulación:} genera eventos y encuestas con periodo, sede, rol, módulo, origen y semilla.
  \item \textbf{Medición:} agrupa observaciones, calcula porcentajes, percentiles, medias y desviación entre sedes.
  \item \textbf{Exportación:} guarda CSV y JSON para que API, tablero e informe consuman la misma evidencia.
\end{enumerate}

La API de medición expone salud, incidentes, métricas, resumen y ejecuciones de simulación. React y Flutter acceden mediante el Gateway. El tablero permite leer valores agregados sin incorporar métricas académicas en las tareas del paciente. Esta separación conserva el carácter realista del HIS: la persona clínica atiende; el paciente gestiona su salud; gerencia observa tendencias.

\begin{table}[H]
\centering
\caption{Relación entre fallo, evidencia y medida}
\begin{tabularx}{\textwidth}{p{3.4cm}YY}
\toprule
Fallo controlado & Evidencia generada & KPI afectado \\
\midrule
Guardado HCE lento & Duración y resultado de \texttt{hce\_save} & M-EFI-01 \\
Agendamiento prolongado & Inicio, pasos y finalización de \texttt{appointment} & M-EFE-01 \\
Factura duplicada & Estado y bandera de duplicidad de \texttt{billing} & M-RIE-01 \\
Dosis incorrecta & Resultado de validación de \texttt{prescription} & M-EFE-02 \\
Imagen lenta & Duración de \texttt{imaging} & M-EFI-02 \\
Teleconsulta interrumpida & Estado de \texttt{teleconsultation} & M-CC-02 \\
\bottomrule
\end{tabularx}
\end{table}

\evidence{La automatización reside en \texttt{analytics}; los datos regenerables en \texttt{data/simulated} y \texttt{data/processed}; y los resúmenes seleccionados en \texttt{evidencias/resultados}.}
